\section{Metodologia}

O aparato experimental foi montado conforme a figura \ref{} e o laser foi alinhado. 
Uma lente convergente foi posicionada em frente ao laser para que os padrões formados possam ser observados em distâncias condizentes com as dimensões do laboratório. 
O espelho $M_1$, em frente ao laser, pode ser movido para frente e para trás com um parafuso micrométrico. 
O outro espelho, $M_2$, possui parafusos de regulagem, que foram usados para centralizar o padrão de interferência no anteparo. 

\begin{table}\label{tab: materials}
\caption{Materiais utilizados nos experimentos.}
\begin{tabular}{c|c|c|c}
	

\end{tabular}
\end{table}

\subsection*{Parte 1: determinação do $\lambda$ do laser}
Utilizando o parafuso micrométrico, foi possível variar a posição de $M_1$, e assim alterar o caminho óptico. 
O espelho foi colocado inicialmente em uma posição $x_0$ correspondente a um mínimo de interferência, que foi anotada. 
O parafuso micrométrico foi girado lentamente até que foram contados $m=60$ mínimos. 
O valor final $x_1$ foi medido e foi calculada a diferença $d = |x_1 - x_0|$. 
A luz percorre esse trajeto indo e voltando, e a alavanca que liga o parafuso ao espelho tem um fator de transmissão de 1:10. 
Assim, a diferença de caminho óptico é $2d/10=d/5$.
Utilizando isso na eq. \eqref{eq: lambda}, obtemos
\begin{equation}\label{eq: lambda(d)}
	\lambda_0 = \frac{1}{5}\left(\frac{d}{m}\right)
\end{equation}
O processo foi repetido um total de 5 vezes.

\subsection*{Parte 2: determinação do $n$ do ar}
O índice de refração do ar não é uma constante absoluta. 
Para gases ideais à baixa pressão, temos que o índice de refração aumenta linearmente com a pressão, e que $p = 0 \implies n = n_0 = 1$. i.e., 
\begin{equation}\label{eq: n(p)}
	n(p) = 1 + \kappa p,
\end{equation}

\noindent onde $\kappa = \partial n / \partial p$. Se a pressão do ar em um trecho de comprimento $L$ em um dos braços do intereferômetro sofrer uma variação de pressão de módulo $\Delta p$, equivalente a uma variação de $\Delta n$ no índice de refração, e isso causar uma variação de $m$ mínimos, teremos 
\begin{align}\label{eq: kappa}
	m &= \frac{2L}{\lambda} |\Delta n| \\[4pt]
	\implies \kappa = \frac{\Delta n}{\Delta p} &= \frac{\lambda}{2L} \frac{m}{\Delta p}.
\end{align}

\noindent Assim, basta montar um gráfico de $m \times \Delta p$ e calcular seu coeficiente angular para calcular o valor de $\kappa$ pela eq. \eqref{eq: kappa}. 

Entre o divisor de feixe e um dos espelhos, foi fixada uma célula para armazenamento de ar, que é ligada por uma mangueira a uma bomba de vácuo manual com manômetro. 
A célula foi posicionada de modo a permitir a passagem do feixe indo e voltando, sem alterar sua trajetória. 
Utilizando a bomba, a pressão dentro da célula foi reduzida até o manômetro indicar $ \Delta p \approx $ 800 mbar abaixo da pressão atmosférica. Por causa de má vedação, o ar lentamente retornava para dentro da célula, fazendo a pressão voltar lentamente à inicial. Foi gravado um vídeo em que é possível ver o padrão alternando entre máximos e mínimos a medida que o ponteiro do manômetro vai de -800 mbar à 0 mbar. Desacelerando o vídeo, foi possível contar os mínimos e construir uma relação entre $\Delta p$ e $\Delta n$. Foram considerados 6 vídeos diferentes.

\subsection*{Parte 3: determinação do $n$ do vidro}
Um fino painel de vidro de foi fixo entre o divisor de feixe e $M_2$ por um parafuso, de modo que foi possível rotacionar o painel em torno de um eixo vertical. 
Inicialmente, o painel foi posto perpendicularmente ao feixe ($\theta_0 = 0^\circ$), posição que corresponde a um mínimo.
Então, o painel lentamente foi girado no sentido anti-horário (aumentando $\theta$), até que foram contados $m$ mínimos de interferência. 

Ao girar o painel, a distância que a luz deve se propagar dentro do vidro aumenta, o que por sua vez aumenta o caminho óptico (pois $n_{vidro} > n_{ar}$).
Para pequenas variações de $\theta$ em torno de $0^{\circ}$, $n=n(\theta, m)$ é dado por
\begin{equation}
	n = \frac{ (2s-m\lambda)(1-\cos\theta)+m^2\lambda^2/4s}{   2s(1-cos\theta)-m\lambda   }
\end{equation}
Foram feitas 9 medidas ao todo, e para cada uma, foi anotada o valor final de $\theta$ e número $m$.